\section{Removing a pseudoeffective error}
\label{sec:conversion}

We prove Theorem~\ref{thm:conversion}. Throughout this section
\(L=-K_S\) carries a fixed smooth semipositive metric, and we choose
effective integral divisors
\begin{equation}\label{eq:effective-sequence}
N_j\sim m_jL-D,\qquad m_1<m_2<\cdots.
\end{equation}
We use additive notation for line bundles. A divisor is pseudoeffective
when its numerical class lies in the closed effective cone; on a smooth projective variety
it is equivalent to admitting a singular Hermitian metric with
semipositive curvature current \cite{DPS2001}. Such classes pull back
under dominant maps between smooth projective varieties and push forward
under birational morphisms. 

The proof builds a resolution \(\pi:V\to S\), a morphism
\(f:V\to Y\), and a line bundle \(B\) on the base, together with
a nonzero map \(f^*B\to\ell\pi^*L\). The map will be used to transfer the
sections constructed on the base to positive multiples of \(L\) on
\(S\). Its normalization will also supply
the bounds needed to construct those base sections.

\subsection{The base and the relative section}

\paragraph{A maximal base.}
Call a dominant rational map \(S\dashrightarrow Y\) \emph{dominated
by \(L\)} if, after resolving it to a morphism \(f:V\to Y\),
\begin{equation}\label{eq:domination}
c\pi^*L-f^*H\quad\text{is pseudoeffective}
\end{equation}
for some positive integer \(c\) and ample Cartier divisor \(H\) on
the integral projective variety \(Y\). The constant map is allowed:
the trivial line on a point is ample, and \(L\) is pseudoeffective.
There is therefore a dominated map of largest possible base dimension.

We may choose its base smooth and its function field relatively
algebraically closed in \(\C(S)\). Indeed, normalize the original
base in its relative algebraic closure in \(\C(S)\), a finite
extension, and then resolve the base and the rational map. The pullback
of the original ample divisor to the new smooth base is big. A large
multiple of this pullback dominates an ample divisor in pseudoeffective
order, which preserves \eqref{eq:domination}. The same argument permits
further birational modifications of the base and higher smooth source
models.

This choice has a useful maximality property. If \(A\) is a line
bundle on \(S\) and \(c'L-A\) is pseudoeffective for some
\(c'>0\), then the ratio of any two nonzero sections of \(A\)
lies in \(\C(Y)\). To prove this, suppose their ratio is nonconstant
and resolve the corresponding pencil. Its moving line is the pullback
of \(\OO_{\PP^1}(1)\), and its fixed divisor is effective, so the
pencil is dominated by \(L\). The joint map of this pencil and
\(S\dashrightarrow Y\) is dominated as well: add their two
inequalities on a common resolution and restrict the ample product
line to the joint image. Normalizing and resolving that image preserves
domination. By maximality its dimension is \(\dim Y\), so the
pencil's function is algebraic over \(\C(Y)\). Relative algebraic
closedness places it in \(\C(Y)\), as asserted. A constant ratio
already has this property. When necessary, increasing \(c'\) to an
integer preserves pseudoeffectivity because \(L\) is pseudoeffective.

Apply this property to the effective divisors in
\begin{equation}\label{eq:interpolation}
(m_j-m_2)N_1+(m_2-m_1)N_j
\sim(m_j-m_1)N_2,\qquad j>2.
\end{equation}
Their common class is \((m_j-m_1)(m_2L-D)\), which is bounded
above by \((m_j-m_1)m_2L\) in pseudoeffective order. The equivalence
function in \eqref{eq:interpolation} thus belongs to \(\C(Y)\).
On any normal model resolving the map to \(Y\), a prime is called
\emph{horizontal} if it dominates \(Y\), and \emph{vertical} otherwise.
A base function has order zero at a horizontal prime. If \(a_j\) is
the coefficient of the pullback of \(N_j\) at such a prime, then
\eqref{eq:interpolation} gives
\[
a_j=a_1+\frac{m_j-m_1}{m_2-m_1}(a_2-a_1).
\]
Since \(a_j\ge0\) along an unbounded sequence, \(a_2\ge a_1\).
Consequently the pullback of \(N_2-N_1\) has effective horizontal part.

If \(Y\) is a point, apply this calculation on \(S\) itself: every
prime is horizontal, and \(N_2-N_1\) is effective. Its class is
\((m_2-m_1)L\), proving the theorem in this case. Henceforth
\(\dim Y>0\). The remaining task is to remove the negative vertical
coefficients by subtracting a divisor from the base.

\paragraph{Normalization along base divisors.}
For that subtraction to see every vertical coefficient, we need each
vertical prime to dominate a base divisor. We construct a normal
intermediate model with this property:
\[
\begin{tikzcd}[column sep=large]
V\arrow[r,"\rho"]\arrow[dr,bend right=15,"f"']&
U\arrow[r,"b"]\arrow[d,"h"]&S\\
&Y&
\end{tikzcd}
\qquad \pi=b\rho,\quad f=h\rho.
\]
Here \(b\) and \(\rho\) are birational, \(V\) and \(Y\) are
smooth projective, and every prime divisor of \(U\) that fails to
dominate \(Y\) dominates a prime divisor of \(Y\).

For this construction, begin with a smooth resolution \(V_0\to Y_0\)
of the maximal map. On a dense open of \(Y_0\), its fibers form a
flat family of subschemes of \(V_0\) with a fixed Hilbert polynomial.
Resolve the induced rational map from \(Y_0\) to the projective
Hilbert scheme, obtaining a smooth projective birational base \(Y\)
\cite[Theorem~3.2, p.~260; p.~265]{Grothendieck1961Hilbert}.
Pull back the universal family. Every fiber of this family has
dimension \(r=\dim S-\dim Y\). The reduced closure \(W\) of the
original flat family is integral and maps birationally to \(V_0\).
Its fibers are closed subschemes of the pulled-back Hilbert fibers,
so have dimension at most \(r\). The normalization \(U\to W\)
is finite, and hence has the same fiber-dimension bound over \(Y\).
A prime divisor \(Q\subset U\) mapping into a subset of codimension
at least two would satisfy
\[
\dim Q\le \dim Y-2+r=\dim S-2,
\]
which is impossible. Resolve \(U\) to obtain \(V\). The dimension
bound is needed on \(U\); the smooth space \(V\) will be used for
metrics and vanishing. The preceding maximality argument preserves
\eqref{eq:domination} throughout these modifications.

On this model, let \(s_0\) be the rational section of
\((m_2-m_1)b^*L\) whose divisor is \(b^*(N_2-N_1)\). Its horizontal
part is effective by the preceding calculation. The base divisor to
be subtracted is now determined one prime at a time.

For each prime divisor \(P\subset Y\), subtract the smallest
normalized vertical order:
\begin{equation}\label{eq:minima}
b_P=\min_{Q\mapsto P}
\frac{\ord_Q(s_0)}{\ord_Q(h^*P)}.
\end{equation}
The minimum runs over the finitely many components of \(h^*P\)
dominating \(P\). This set is nonempty and its denominators are
positive. Only finitely many \(b_P\) are nonzero, because
\(\divisor(s_0)\) has finite support. Every vertical prime of \(U\)
occurs in one of these minima. Together with horizontal effectivity,
this proves
\[
\divisor(s_0)-h^*\Bigl(\sum_P b_PP\Bigr)\ge0.
\]
Choose a positive integer \(a\) clearing all denominators and set
\[
\ell=a(m_2-m_1),\qquad B=a\sum_Pb_PP.
\]
Smoothness of \(Y\) makes \(B\) Cartier. On the normal variety
\(U\), the resulting rational section of \(\ell b^*L-h^*B\)
has nonnegative order at every prime, so is regular. Pulling it back
to \(V\) gives
\begin{equation}\label{eq:sigma}
\sigma:\OO_V(f^*B)\longrightarrow\OO_V(\ell\pi^*L).
\end{equation}
For every base prime \(P\), some component of \(f^*P\) dominating
\(P\) has order zero in \(\sigma\): take the strict transform of
a prime attaining \eqref{eq:minima}. Its strict transform defines the
same divisorial valuation, since \(U\) is normal and hence regular
at its generic point. The new exceptional divisors on \(V\) acquire
no poles because the pulled-back section is regular. This is the
normalization we will use near singular fibers; equidimensionality of
\(V\) is unnecessary.

\paragraph{Rank one of the adjoint direct images.}
We now have the map that will transfer sections from the base. To
construct a semipositive metric on its source line \(B\), we will
use adjoint integral metrics. Their direct images must first be shown
to have rank one. The anticanonical identity enters at this point. Put
\begin{equation}\label{eq:discrepancy}
E=K_V-\pi^*K_S=K_V+\pi^*L.
\end{equation}
The divisor \(E\) is effective and \(\pi\)-exceptional; denote its
canonical section by \(e\). For every integer \(k\ge0\), we claim
that
\begin{equation}\label{eq:rank-one}
\cG_k=f_*\OO_V(E+k\ell\pi^*L)
\quad\text{has generic rank one.}
\end{equation}
The section \(e\sigma^k\), after trivializing \(B\) over the
generic point, shows that the rank is at least one, also when \(k=0\).
If it were larger, Serre's global generation theorem would give, for
some positive \(t\), two sections of
\(E+k\ell\pi^*L+tf^*H\) independent over \(\C(Y)\). Their ratio
would lie outside that field.

Push their zero divisors forward to \(S\). They become two effective
divisors in one Cartier class
\[
A\sim k\ell L+t\pi_*f^*H,
\]
with the same rational-function ratio. Indeed, choose a Cartier divisor
representing the upstairs line and express both sections as rational
functions; birational pushforward preserves their principal divisors,
and \(\pi_*E=0\). Smoothness makes the common downstairs divisor
Cartier. Pushing forward \eqref{eq:domination} gives
\[
(k\ell+tc)L-A\quad\text{pseudoeffective}.
\]
The maximality property then places the ratio in \(\C(Y)\), a
contradiction. The coefficient \(k\ell+tc\) is positive even at
\(k=0\), proving \eqref{eq:rank-one} in its full stated range.

\subsection{The bounded base metric}

We have constructed \(B\) and \(\sigma\); the next task is a
semipositive metric on \(B\) with locally bounded weights. For a local
frame \(\beta\), our weight convention is \(-\log|\beta|^2\),
so semipositive curvature means that the weights are plurisubharmonic.
We obtain them as limits of adjoint integral metrics. We use
Berndtsson's direct-image theorem in the following
form \cite[Theorem~1.2]{Berndtsson2009}: for a proper holomorphic
submersion with K\"ahler total space and a smoothly semipositive line
bundle \(A\), the natural \(L^2\) metric on
\(f_*(K_{V/Y}+A)\) has semipositive curvature wherever the adjoint
sections form a vector bundle. The coefficient metric may be
semipositive; strict positivity along the fibers is not required.

Let \(h_L\) denote the pulled-back smooth metric on \(\pi^*L\).
Choose a dense Zariski open \(Y^\circ\) on which \(f\) is smooth
and \(e\) and \(\sigma\) are not identically zero on any fiber.
The fibers there are connected: relative algebraic closedness of
\(\C(Y)\) in \(\C(S)\) makes the finite part of the Stein
factorization trivial. For a local frame \(\beta\) of \(B\), set
\(\sigma_\beta=\sigma(f^*\beta)\) and
\begin{equation}\label{eq:max-weight}
\varphi_\beta(y)=-\log\max_{V_y}
|\sigma_\beta|_{h_L^\ell}^{\,2},\qquad y\in Y^\circ.
\end{equation}
The fiberwise maximum is positive and continuous on this open set.

To see that \(\varphi_\beta\) is plurisubharmonic, take a local
frame \(\tau\) of \(K_Y\). Equation~\eqref{eq:discrepancy} gives
\begin{equation}\label{eq:adjoint}
K_{V/Y}+(1+k\ell)\pi^*L
=E+k\ell\pi^*L-f^*K_Y.
\end{equation}
Thus \(e\sigma_\beta^k/f^*\tau\) is an adjoint section. By
\eqref{eq:rank-one} and generic base change, it spans the adjoint
direct image on a dense open subset of \(Y^\circ\). For every
integer \(k\ge1\), Berndtsson's theorem makes
\begin{equation}\label{eq:moment-weights}
\varphi_{\beta,k}(y)
=-\frac1k\log\int_{V_y}
|\sigma_\beta|_{h_L^\ell}^{\,2k}
\left|e/f^*\tau\right|_{h_L}^{\,2}
\end{equation}
plurisubharmonic on that dense open. The integral is smooth and
strictly positive throughout \(Y^\circ\); its curvature inequality
therefore holds on all of \(Y^\circ\) by continuity.

The final factor in the integral is the smooth nonnegative measure
obtained from a \(\pi^*L\)-valued relative canonical form. It has
positive mass and full support on each fiber, because a nonzero
holomorphic section on a connected smooth fiber cannot vanish on an
open subset. Denote this measure by \(\mu_y\), its mass by
\(M(y)>0\), and put
\[
q_\beta=|\sigma_\beta|_{h_L^\ell}^{\,2},\qquad
A_k(y)=\left(\frac1{M(y)}\int_{V_y}q_\beta^k\,d\mu_y\right)^{1/k}.
\]
H\"older's inequality for the probability measure \(M(y)^{-1}\mu_y\)
shows that \(A_k(y)\) increases with \(k\). Full support gives
\(A_k(y)\to\max_{V_y}q_\beta\). The moments and their limit are
continuous, so Dini's theorem makes this convergence uniform on every
compact subset of a coordinate open in \(Y^\circ\). The limit is
positive there. Since
\[
\varphi_{\beta,k}=-\log A_k-\frac1k\log M,
\]
the weights in \eqref{eq:moment-weights} converge locally uniformly to
\(\varphi_\beta\). Their limit is plurisubharmonic.

It remains to establish local bounds near the omitted fibers.
Properness of \(f\) and smoothness of the upstairs metric bound
\(q_\beta\) above over every relatively compact base neighborhood.
Equation~\eqref{eq:max-weight} therefore gives a local lower bound
for \(\varphi_\beta\), even as one approaches the boundary.

For the upper bound, let \(P\) be a divisorial component of
\(Y\setminus Y^\circ\). The normalization of \(\sigma\) supplies
a component \(Q\) above \(P\) on which \(\sigma\) is generically
nonzero. Choose a general point of \(Q\) where \(P,Q\) are smooth,
\(Q\to P\) is submersive, no other component of \(f^*P\) passes
through the point, and \(\sigma\ne0\). In suitable local coordinates, a
transverse coordinate to \(P\) pulls back to a unit times \(z_1^r\),
with \(r>0\), while the other base coordinates pull back to independent
coordinates along \(Q\). Taking a local root absorbs the unit.
This local form shows that a small source neighborhood maps onto a
base neighborhood. On a still smaller source neighborhood,
\(q_\beta\) has a positive lower bound. Every sufficiently nearby
fiber meets this neighborhood, so its maximum has the same lower
bound. This gives an upper bound for \(\varphi_\beta\) near a
general point of \(P\).

The removable-singularity theorem for plurisubharmonic functions now
extends \(\varphi_\beta\) across dense open subsets of all boundary
divisors. These opens can be taken Zariski open: nonvanishing and the
required differential ranks hold on algebraic opens of \(Q\), whose
dominant constructible images contain dense Zariski opens of \(P\).
The set still omitted is therefore contained in a closed analytic set
\(A\subset Y\) of codimension at least two.

We recall why no upper-bound obstruction remains at \(A\). Near a
point of \(A\), choose a small affine complex disk centered there whose
boundary misses \(A\), and a compact neighborhood of that boundary disjoint
from \(A\). The plurisubharmonic function has a common upper bound
on this neighborhood. For every nearby \(y\notin A\), choose a disk
of the same radius centered at \(y\), with direction sufficiently
close to the original one that its boundary stays in this fixed
neighborhood. Its direction can also be chosen so that the whole
disk misses \(A\):
the radial image of \(A\), viewed from \(y\), in the space of complex line directions
has real dimension at most \(2\dim Y-4\), smaller than the direction
space dimension \(2\dim Y-2\). The submean inequality gives
a common upper bound at all such \(y\). Removable singularities
therefore extend the function across \(A\) as well. The upper-limit
extension preserves the lower bound already supplied by properness.

We have obtained locally bounded plurisubharmonic weights on all of
\(Y\). Under \(\beta'=g\beta\), formula~\eqref{eq:max-weight}
gives
\[
\varphi_{\beta'}=\varphi_\beta-\log|g|^2.
\]
The same relation holds after extension, so the weights define a
semipositive singular Hermitian metric \(h_B\) on \(B\), with
locally bounded weights.

\subsection{Vanishing at every nonnegative exponent}

The bounded metric on \(B\) is now available. To produce sections,
we combine it with a metric on \(\pi^*L\) that is positive in base
directions and has no multiplier-ideal loss. Recall that \(\cJ(h)\)
consists locally of holomorphic functions \(g\) for which
\(|g|^2e^{-\varphi_h}\) is integrable.

The domination \eqref{eq:domination} gives a singular metric on
\(\pi^*L\) whose curvature is positive in base directions. More
precisely, combine a semipositive singular metric on
\(c\pi^*L-f^*H\) with a positive smooth metric on \(H\), and take
the \(c\)-th root. The resulting metric \(h_1\) satisfies
\[
\Theta_{h_1}(\pi^*L)\ge c^{-1}f^*\omega_H.
\]
Mix its local weights with those of the smooth semipositive metric
\(h_L\): for some \(0<\delta<1\), put
\(h=h_L^{1-\delta}h_1^\delta\). The local exponential-integrability consequence of Skoda's theorem
\cite[Proposition~7.1, p.~389]{Skoda1972} supplies, near each point,
a sufficiently small positive exponent for which the negative
exponential of the singular weight is integrable. We use its
arbitrary-plurisubharmonic-weight formulation in
\cite[Definition~0.1 and 1.4(8)]{DemaillyKollar2001}:
for a nontrivial psh weight \(\varphi\), the local integrability
threshold of \(e^{-2\gamma\varphi}\) is positive. This formulation
also includes complex dimension one. In our squared-norm convention,
take \(\gamma=\delta/2\). A finite coordinate
cover of the compact space \(V\) gives one common \(\delta>0\).
The smooth part changes integrability by bounded factors. Hence
\begin{equation}\label{eq:small-metric}
\Theta_h(\pi^*L)\ge\epsilon f^*\omega_H,
\qquad \cJ(h)=\OO_V,
\qquad \epsilon=\delta/c>0.
\end{equation}

For each integer \(k\ge0\), tensor this metric with \(f^*h_B^k\).
Its curvature is still at least \(\epsilon f^*\omega_H\), and its
multiplier ideal is still trivial: for every fixed \(k\), the added
weight is locally bounded. Thus the same choice of \(h\) works at
all the required exponents.

We use Fujino's Koll\'ar--Nadel vanishing theorem in the following
form \cite[Theorem~1.3]{Fujino2018}. If \(f:V\to Y\) is a
surjective morphism from a compact K\"ahler manifold to a projective
variety and a line bundle \(A\) has a singular metric \(h_A\) with
\(\Theta_{h_A}(A)\ge\epsilon f^*\omega_H\), then
\[
H^i\bigl(Y,R^qf_*(\OO_V(K_V+A)\otimes\cJ(h_A))\bigr)=0
\qquad (i>0,\ q\ge0).
\]
Fujino derives this form from the injectivity and multiplier-ideal
Bertini theorems of Fujino--Matsumura
\cite[Theorems~1.1--1.3]{Fujino2018}. The metric need not have analytic
singularities. Apply the theorem to
\(A=\pi^*L+kf^*B\), with the metric just constructed, and take \(q=0\). Since
\(K_V+\pi^*L=E\), the projection formula gives
\begin{equation}\label{eq:all-vanishing}
H^i\bigl(Y,\cG\otimes\OO_Y(kB)\bigr)=0
\qquad(i>0,\ k\ge0),
\qquad \cG=f_*\OO_V(E).
\end{equation}

Riemann--Roch on the smooth projective base makes
\[
Q(k)=\chi\bigl(Y,\cG\otimes\OO_Y(kB)\bigr)
=\int_Y\ch(\cG)e^{k c_1(B)}\td(Y)
\]
a polynomial in \(k\), whether or not \(B\) is ample. By
\eqref{eq:all-vanishing}, this polynomial equals
\(h^0(Y,\cG\otimes\OO_Y(kB))\) at every nonnegative integer.
At zero the canonical section \(e\) gives a nonzero section, so
\(Q(0)>0\). This is why vanishing at \(k=0\) is indispensable:
without it, the existence of \(e\) would not determine the Euler
characteristic. The polynomial is nonzero, and hence cannot vanish
at every positive integer. For some \(k>0\) there is a nonzero
section of \(E+kf^*B\) on \(V\).

Multiply this section by \(\sigma^k\) from \eqref{eq:sigma}. The
result is a nonzero section of \(E+k\ell\pi^*L\). Finally,
\(\pi_*\OO_V(E)=\OO_S\): a rational function with possible poles
only on the effective exceptional divisor \(E\) is regular away
from its codimension-two image, and extends across that image by
normality of \(S\). The projection formula therefore sends our
section to a nonzero section of \(k\ell L\), completing the proof
of Theorem~\ref{thm:conversion}.

The two appearances of \(L=-K_S\) explain the strength and the scope
of the criterion. The identity \(K_V+\pi^*L=E\) supplies the
adjoint section used in the integral metric and turns adjoint
vanishing into the fixed sheaf \(f_*\OO_V(E)\) at exponent zero.
These steps connect the positive base metric to an Euler polynomial
with a known nonzero value.
